% GENERATED FILE. Edit canonical lessons, then run npm run generate:books.
% canonical-source-hash: fnv1a64:d8bc1786c74e6bfb
% canonical-lessons: LA-C255-teneo, LA-C255-rumpo, LA-C255-calefacio, LA-C255-impleo, LA-C255-metior

\chapter{To hold, To burst, To warm, To fill, To measure}
\label{ch:la-to-hold-to-burst-to-warm-to-fill-to-measure}
\hlchaptermodality{255}

\begin{chapteropening}
\textbf{By the end of this chapter:} \emph{I can say to hold, to burst, to warm, to fill and to measure.}

Complete the last lesson of chapter 255: I can say to hold, to burst, to warm, to fill and to measure.
\end{chapteropening}

\section[teneō, tenēre]{teneō, tenēre — to hold}

\label{lesson:LA-C255-teneo}



\begin{quote}
Before the new one: say the Latin for a tent, then the Latin for a flag, a banner.
\end{quote}

\subsection*{You'll want to know: teneō, tenēre}
\textbf{teneō, tenēre} — \textquotedblleft{}to hold\textquotedblright{}.

\begin{cousinweb}
\textbf{teneō} is inherited from Proto-Italic \emph{*tenēō}, \textquotedblleft{}to hold\textquotedblright{} (EtymDB 2.1), usually traced to Proto-Indo-European \emph{*ten-}, \textquotedblleft{}to stretch\textquotedblright{} (Lewis \& Short compare Greek \emph{teinō}, \textquotedblleft{}to stretch\textquotedblright{}). English \emph{tenant} comes from it.
\end{cousinweb}

\begin{grammarlens}[title={What you've built}]
One more word for this chapter.
\end{grammarlens}

\subsection*{Your turn}
\begin{itemize}
\raggedright
  \item \emph{Say it:} \emph{teneō, tenēre}
  \item \emph{Say it:} \emph{teneō, tenēre}, once more
  \item \emph{Say it:} say \emph{teneō, tenēre}
  \item \emph{Recall:} say the Latin for a tent, then the Latin for a flag, a banner, then say \emph{teneō, tenēre} again
\end{itemize}

\begin{tcolorbox}[breakable,colback=teal!4,colframe=teal!35!black,title={Before you move on}]
What does teneō, tenēre mean? (\textquotedblleft{}To hold\textquotedblright{}.) Say it once more.
\end{tcolorbox}
\section[rumpō, rumpere]{rumpō, rumpere — to burst, to break apart}

\label{lesson:LA-C255-rumpo}



\begin{quote}
Before the new one: say the Latin for a flag, a banner, then the Latin for to hold.
\end{quote}

\subsection*{You'll want to know: rumpō, rumpere}
\textbf{rumpō, rumpere} — \textquotedblleft{}to burst, to break apart\textquotedblright{}.

\begin{cousinweb}
\textbf{rumpō} is inherited, usually traced to a Proto-Indo-European present \emph{*Hrunepti}, \textquotedblleft{}breaks\textquotedblright{} (EtymDB 2.1), with an n inside the root. English \emph{rupture} comes from it.
\end{cousinweb}

\begin{grammarlens}[title={What you've built}]
One more word for this chapter.
\end{grammarlens}

\subsection*{Your turn}
\begin{itemize}
\raggedright
  \item \emph{Say it:} \emph{rumpō, rumpere}
  \item \emph{Say it:} \emph{rumpō, rumpere}, once more
  \item \emph{Say it:} say \emph{rumpō, rumpere}
  \item \emph{Recall:} say the Latin for a flag, a banner, then the Latin for to hold, then say \emph{rumpō, rumpere} again
\end{itemize}

\begin{tcolorbox}[breakable,colback=teal!4,colframe=teal!35!black,title={Before you move on}]
What does rumpō, rumpere mean? (\textquotedblleft{}To burst, to break apart\textquotedblright{}.) Say it once more.
\end{tcolorbox}
\section[calefaciō, calefacere]{calefaciō, calefacere — to warm, to heat}

\label{lesson:LA-C255-calefacio}



\begin{quote}
Before the new one: say the Latin for to hold, then the Latin for to burst.
\end{quote}

\subsection*{You'll want to know: calefaciō, calefacere}
\textbf{calefaciō, calefacere} — \textquotedblleft{}to warm, to heat\textquotedblright{}.

\begin{cousinweb}
\textbf{calefaciō} is built from \emph{faciō}, \textquotedblleft{}to make\textquotedblright{} (EtymDB 2.1), with the stem of \emph{caleō}, \textquotedblleft{}to be warm\textquotedblright{}: to make warm.
\end{cousinweb}

\begin{grammarlens}[title={What you've built}]
One more word for this chapter.
\end{grammarlens}

\subsection*{Your turn}
\begin{itemize}
\raggedright
  \item \emph{Say it:} \emph{calefaciō, calefacere}
  \item \emph{Say it:} \emph{calefaciō, calefacere}, once more
  \item \emph{Say it:} say \emph{calefaciō, calefacere}
  \item \emph{Recall:} say the Latin for to hold, then the Latin for to burst, then say \emph{calefaciō, calefacere} again
\end{itemize}

\begin{tcolorbox}[breakable,colback=teal!4,colframe=teal!35!black,title={Before you move on}]
What does calefaciō, calefacere mean? (\textquotedblleft{}To warm, to heat\textquotedblright{}.) Say it once more.
\end{tcolorbox}
\section[impleō, implēre]{impleō, implēre — to fill}

\label{lesson:LA-C255-impleo}



\begin{quote}
Before the new one: say the Latin for to burst, then the Latin for to warm.
\end{quote}

\subsection*{You'll want to know: impleō, implēre}
\textbf{impleō, implēre} — \textquotedblleft{}to fill\textquotedblright{}.

\begin{cousinweb}
\textbf{impleō} is \emph{in}, \textquotedblleft{}in\textquotedblright{}, with \emph{pleō}, \textquotedblleft{}to fill\textquotedblright{} (EtymDB 2.1), a verb found only in compounds. \emph{pleō} is usually traced to Proto-Indo-European \emph{*pleh1-}, \textquotedblleft{}to fill\textquotedblright{}; English \emph{full} is a cousin.
\end{cousinweb}

\begin{grammarlens}[title={What you've built}]
One more word for this chapter.
\end{grammarlens}

\subsection*{Your turn}
\begin{itemize}
\raggedright
  \item \emph{Say it:} \emph{impleō, implēre}
  \item \emph{Say it:} \emph{impleō, implēre}, once more
  \item \emph{Say it:} say \emph{impleō, implēre}
  \item \emph{Recall:} say the Latin for to burst, then the Latin for to warm, then say \emph{impleō, implēre} again
\end{itemize}

\begin{tcolorbox}[breakable,colback=teal!4,colframe=teal!35!black,title={Before you move on}]
What does impleō, implēre mean? (\textquotedblleft{}To fill\textquotedblright{}.) Say it once more.
\end{tcolorbox}
\section[mētior, mētīrī]{mētior, mētīrī — to measure}

\label{lesson:LA-C255-metior}



\begin{quote}
Before the new one: say the Latin for to warm, then the Latin for to fill.
\end{quote}

\subsection*{You'll want to know: mētior, mētīrī}
\textbf{mētior, mētīrī} — \textquotedblleft{}to measure\textquotedblright{}.

\begin{cousinweb}
\textbf{mētior} is inherited, usually traced to Proto-Indo-European \emph{*meh1-}, \textquotedblleft{}to measure\textquotedblright{} (EtymDB 2.1); Greek \emph{metron}, \textquotedblleft{}a measure\textquotedblright{}, is a cousin (Lewis \& Short).
\end{cousinweb}

\begin{grammarlens}[title={What you've built}]
One more word for this chapter.
\end{grammarlens}

\subsection*{Your turn}
\begin{itemize}
\raggedright
  \item \emph{Say it:} \emph{mētior, mētīrī}
  \item \emph{Say it:} \emph{mētior, mētīrī}, once more
  \item \emph{Say it:} say \emph{mētior, mētīrī}
  \item \emph{Recall:} say the Latin for to warm, then the Latin for to fill, then say \emph{mētior, mētīrī} again
\end{itemize}

\begin{tcolorbox}[breakable,colback=teal!4,colframe=teal!35!black,title={Before you move on}]
What does mētior, mētīrī mean? (\textquotedblleft{}To measure\textquotedblright{}.) Say it once more.
\end{tcolorbox}
